%% gaze_paper.tex — IEEE conference paper
%%
%% A Quantum Conditional Selection Rule for Webcam Gaze Typing,
%% and Canthus-Relative Iris Normalisation
%%
%% Built on IEEEtran.cls v1.8b. Note: the supplied skeleton was bare_jrnl.tex
%% (the *journal* template, which carries \markboth headers and IEEEbiography
%% blocks). This file uses the [conference] class option instead, which is what
%% a conference submission requires; the biography environments are therefore
%% omitted. Switch the first line back to [journal] if a transactions
%% submission is intended.
%%
%% Compile:  pdflatex gaze_paper && bibtex gaze_paper && pdflatex gaze_paper x2
%% (the bibliography is inlined as thebibliography, so bibtex is not required)

\documentclass[conference]{IEEEtran}

\usepackage{cite}
\usepackage{amsmath}
\usepackage{amssymb}
\usepackage{graphicx}
\usepackage{url}
\usepackage{booktabs}
\usepackage{multirow}
\usepackage{array}
\usepackage{xcolor}

% Placeholder frames for figures that are not yet plotted. Replace each
% \figholder{...} with \includegraphics[width=\linewidth]{yourfile} when the
% plots are exported. The reserved heights are deliberate: reviewers respond
% badly to a paper whose figures are squeezed, and the space is easier to fill
% than to find later.
\newcommand{\figholder}[2]{%
  \fboxsep=0pt\fcolorbox{black!25}{black!3}{%
    \parbox[c][#1][c]{\linewidth}{\centering\footnotesize\color{black!55}#2}}}

% IEEEtran declares no theorem environments; the two propositions in
% Section V need this.
\newtheorem{proposition}{Proposition}

\interdisplaylinepenalty=2500

\hyphenation{nor-mal-isa-tion cal-i-bra-tion land-mark}

\begin{document}

\title{A Quantum Conditional Selection Rule for\\ Low-Cost Webcam Gaze Typing}

\author{
\IEEEauthorblockN{First Author}
\IEEEauthorblockA{Department of Computer Science and Engineering\\
Institution Name\\
City, Country\\
Email: first.author@institution.edu}
\and
\IEEEauthorblockN{Second Author}
\IEEEauthorblockA{Department of Computer Science and Engineering\\
Institution Name\\
City, Country\\
Email: second.author@institution.edu}
}

\maketitle

\begin{abstract}
Text entry still assumes working hands and a physical keyboard, excluding users
with motor impairments and failing in hands-occupied or sterile settings. We
present an eye-gaze typing system that runs entirely on a commodity webcam, with
no infrared hardware and no client build step. Gaze is estimated by mapping
MediaPipe iris landmarks through a per-user ridge-regression calibration, and
characters are committed by dwell selection over either a per-letter keyboard or
large, gaze-friendly key clusters. Clustering makes targets easy to hit but
renders the input ambiguous: several letters share one region, so distinct words
collapse onto the same cluster code. To resolve that ambiguity accurately and in
real time we introduce a \emph{quantum conditional selection rule}. Candidate
targets form a superposition whose amplitudes carry the gaze likelihood; a
language-conditioned controlled unitary fuses linguistic context through both
amplitude and phase; and a Born-rule measurement commits the most probable
candidate. Unlike classical Bayesian fusion, which multiplies marginals, the
phase term admits interference between context-coherent hypotheses. We prove the
rule is strictly more expressive --- it recovers log-linear fusion exactly at
zero phase, and no classical mixture reproduces its cross-terms --- and a
consistency-preserving oracle guarantees that every decoded output re-encodes to
exactly the gazed clusters, never an off-target word. The rule is realised as a
lightweight classical state-vector simulation over a five-qubit register, well
inside the webcam frame budget.
We then measure it, and report that the added expressiveness does not pay. At
matched prompt rate the interference term is worth a median $-0.03$ percentage
points across five noise levels, with no significant win or loss; amplitude
amplification lengthens mean dwell from $2.63$ to $2.80$ frames and costs two
points of commit accuracy. We give the mechanism rather than the shrug:
coherence is \emph{anti}-correlated with correctness, because the correct letter
after an English prefix carries the richest hypothesis support and therefore the
widest spread of phases. The accuracy instead comes from the front end.
Replacing eyelid-relative normalisation of the iris offset with a
canthus-relative one reduces held-out error from $5.97^\circ$ to $4.37^\circ$, a
$26.5\,\%$ reduction that wins on every recorded session, subsumes explicit head
pose, and is invariant to head roll by construction to $10^{-15}$. All results
come from three recorded sessions of a single participant; we state this
prominently rather than in closing.
\end{abstract}

\begin{IEEEkeywords}
Quantum conditional selection rule; quantum-inspired inference;
interference-based gaze--language fusion; amplitude amplification;
cluster-constrained decoding; eye-gaze typing; webcam gaze estimation;
dwell-based selection; assistive text entry; facial landmarks; calibration.
\end{IEEEkeywords}

\IEEEpeerreviewmaketitle


\section{Introduction}

\IEEEPARstart{E}{ye} typing offers people with severe motor impairment a channel
for written communication that requires no limb movement. Commercial infrared
eye trackers achieve roughly $0.25$--$0.5^\circ$ of gaze error but cost
substantially more than the computers they attach to, so a large body of work
has pursued the same capability using only a consumer webcam.

A webcam gaze keyboard faces two problems in series, and they are not the same
kind of problem. The first is estimation: turning a low-resolution, uncalibrated
image of a face into a screen coordinate. The second is decision: turning a noisy
coordinate and a language context into a committed character. Most of the
literature addresses the first and treats the second as a hit test.

This paper addresses both, and finds the balance between them counter-intuitive.

For the decision problem we introduce a \emph{quantum conditional selection
rule}. The deployed system already fuses gaze and language the conventional way,
by multiplying a gaze posterior with a language prior. That product is a product
of \emph{marginals}, and a marginal has already discarded the structure that
produced it: whether a letter's linguistic support came from one confident word
or from fifty diffuse ones. The selection rule keeps the hypotheses distinct,
places them in superposition with amplitudes carrying the gaze likelihood,
applies a language-conditioned unitary that attaches both a magnitude and a
\emph{phase} to each, and reads the result with the Born rule. The phase channel
admits interference, which no product of marginals can express.

That is a genuine generalisation and we prove it. It is also, we find, an
empirically empty one, and the more useful contribution of this paper is the
explanation of why. The interference term is systematically mis-signed for this
task: the correct letter after an English prefix is usually the one with the
richest vocabulary support, so it carries the most hypotheses, the widest spread
of phases, and hence the \emph{lowest} coherence. Phase derived from hypothesis
likelihood penalises exactly the evidence it should reward, and no retuning of
the phase function repairs the sign.

For the estimation problem the literature has converged on an evaluation habit
that, we argue, obscures progress. Systems are reported as \emph{classifiers}:
the metric is the proportion of frames in which the gaze direction is assigned to
the correct region, where the regions number between two and nine. Reported
figures of $85$--$93\,\%$ are common. This is a legitimate measurement of a
legitimate design --- most such systems select keys by blinking or by scanning,
so a coarse direction is all they require --- but it is not comparable to the
continuous angular error that the wider gaze literature reports, and it cannot
tell a reader whether a user could point at a key.

We take the harder formulation: a continuous screen coordinate
(Section~\ref{sec:frontend}), evaluated on targets never used for fitting under
the protocol of Section~\ref{sec:protocol}, and reported in degrees of visual
angle. Working in
that regime surfaced a defect in a step the literature treats as settled, and it
is there --- not in the decision rule --- that the accuracy turned out to be.

\subsection{Contributions}

\begin{enumerate}
  \item \textbf{A quantum conditional selection rule} for gaze--language fusion:
  candidates in superposition with gaze-derived amplitudes, a
  language-conditioned controlled unitary carrying context in amplitude and
  phase, a consistency-preserving oracle, and Born-rule commitment, realised as
  classical state-vector simulation over a five-qubit register.

  \item \textbf{A containment proposition.} The rule reduces to the deployed
  log-linear fusion exactly, to floating point, at zero phase
  (Proposition~\ref{prop:contain}), and its cross-terms are unreachable by any
  classical mixture of the same hypotheses (Proposition~\ref{prop:express}).
  Containment is what makes the head-to-head comparison honest rather than
  rigged.

  \item \textbf{A measured null with a mechanism.} Interference is worth a
  median $-0.03$ percentage points at matched prompt rate and amplification
  lengthens dwell. We trace this to a sign: the measured coherence separation
  between correct and incorrect letters is negative at every phase scale and
  worsens monotonically as the scale grows.

  \item \textbf{Canthus-relative normalisation.} Replacing eyelid-relative
  normalisation of the iris offset with a canthus-relative one reduces held-out
  error by $26.5\,\%$ ($5.97^\circ \rightarrow 4.37^\circ$), winning on every
  recorded session. The mechanism is anatomical: the eyelid moves independently
  of gaze, the canthi do not. To our knowledge this normalisation appears in
  none of the surveyed works.

  \item \textbf{An explanation for a reported but undiagnosed asymmetry.} One
  surveyed system reports $93\,\%$ horizontal against $76\,\%$ vertical accuracy
  and does not account for the gap. Eyelid contamination predicts exactly this
  asymmetry, since only the vertical channel is normalised by the aperture.

  \item \textbf{Roll invariance by construction.} Projecting the offset onto the
  canthus axis and its perpendicular removes head-roll dependence without
  estimating the roll angle. We verify invariance numerically to
  $\sim$$10^{-15}$ across $0$--$40^\circ$.

  \item \textbf{An error budget that inverts with head pose.} Per-frame noise is
  $5.3\,\%$ of the error with a still head and $57.9\,\%$ past $10^\circ$ of
  roll. We show the instability originates in the landmark detector rather than
  in the regressor, whose amplification of it stays flat.

  \item \textbf{Reported negative results,} including quantum-kernel regression,
  explicit head-pose features, an iris-radius scale cue, eye-region de-rotation
  and discrete hierarchical selection --- each implemented, measured, rejected,
  and recorded so the next group need not repay the cost.
\end{enumerate}


\section{Related Work}

Table~\ref{tab:survey} summarises ten recent works on webcam gaze keyboards and
related interfaces. We group them by how the iris is localised, since that
choice determines what the rest of the system can do.

\textbf{Threshold and projection methods.} Saraswati \emph{et al.}~\cite{ref08}
detect the face with a Haar cascade and locate the iris by binary thresholding
followed by integral projection. The approach is inexpensive but the threshold
must be retuned for each lighting condition --- the paper tabulates separate
values for $8$, $30$ and $81$\,lx --- and detection fails beyond $50$\,cm.

\textbf{Geometric circle fitting.} Cabatuan \emph{et al.}~\cite{ref09} treat the
iris as a circle-detection problem, applying Sobel edge detection, a thinning
step and a Circular Hough Transform, and implement the pipeline on an FPGA in
Verilog. It is the only surveyed work to target dedicated hardware.

\textbf{Landmark regression with aspect ratios.} The largest group
\cite{ref02,ref03,ref05,ref07} uses the 68-point \emph{dlib} shape predictor and
derives an Eye Aspect Ratio (EAR) or a horizontal/vertical gaze ratio. A
property of this family deserves emphasis: the 68-point model contains
\emph{no iris landmarks at all}. It locates the eye corners and lids, and
everything concerning the pupil is inferred from those. This is why all four of
these systems select keys by \emph{blinking} or scanning rather than by
pointing --- the representation cannot support a continuous coordinate.

\textbf{Dedicated iris landmarks.} Rafid \emph{et al.}~\cite{ref04} and the
virtual-mouse system of~\cite{ref06} use MediaPipe Face Mesh, which supplies a
five-point iris ring per eye in addition to 468 face landmarks. This yields the
iris centre directly and is the prerequisite for continuous regression. Only two
of the ten surveyed works adopt it, both dated 2023 or later.

\textbf{End-to-end learning.} Waideman and Aquino Junior~\cite{ref01} regress
gaze from pixels with a four-stream CNN and are the only surveyed work to model
head pose, supplying head position as a fourth network channel. Abhaya
\emph{et al.}~\cite{ref03} classify nine gaze directions with a CNN.

\textbf{Selection strategy.} Zhou \emph{et al.}~\cite{ref09b} address the Midas
touch problem with an adaptive dwell time that shortens as the user's input
speed rises, using a commercial SMI eye tracker rather than a webcam.

\textbf{Ambiguous input and language models.} Enlarging targets to make them
easy to hit, then recovering the intended word with a language model, is the
established trade in text entry: SHARK\textsuperscript{2}~\cite{ref17} decodes
whole-word gestures against a lexicon, and the same bargain underlies our
cluster layout. Our addition is the constraint direction. Where a conventional
decoder scores hypotheses and may return a word the input does not support, the
consistency oracle of Section~\ref{sec:rule} admits only hypotheses that
re-encode to exactly the gazed sequence.

\subsection{Positioning}

Two observations follow from Table~\ref{tab:survey} and motivate this work.

First, \textbf{no surveyed work reports gaze error in degrees}. The reported
percentages measure region classification, often over as few as two or three
zones where chance alone is $33$--$50\,\%$. Our results therefore cannot be
ranked against them, and we do not attempt to; we state this explicitly because
a table of mixed metrics invites exactly that error.

Second, \textbf{head pose is almost entirely absent}. One work models it
implicitly~\cite{ref01}; one \emph{prohibits} head movement and reports the
resulting discomfort as a limitation~\cite{ref02}; one observes that head
shaking causes error and does not address it~\cite{ref09b}. No surveyed work
performs geometric head-pose estimation. Section~\ref{sec:results} offers a
partial explanation for why that omission has cost less than might be expected.

% ---------------------------------------------------------------------------
\begin{table*}[!t]
\renewcommand{\arraystretch}{1.25}
\caption{Survey of Ten Recent Webcam Gaze-Keyboard Systems. The Metric Column is
the Point of the Table: Reported Figures Measure Region \emph{Classification},
Not Continuous Gaze Error, and Are Therefore Not Comparable With This Work.}
\label{tab:survey}
\centering
\footnotesize
\begin{tabular}{@{}p{2.6cm}p{1.9cm}p{3.2cm}p{3.0cm}p{2.0cm}p{2.4cm}@{}}
\toprule
\textbf{Work} & \textbf{Face} & \textbf{Iris localisation} &
\textbf{Selection} & \textbf{Head pose} & \textbf{Reported} \\
\midrule
Waideman \emph{et al.}~\cite{ref01} & MMOD-CNN & 4-stream CNN (face, both eyes, head) & CNN regression & \textbf{Yes}, 4th channel & 93\,\% L/R, \textbf{76\,\% U/D} \\
Islam \emph{et al.}~\cite{ref02} & HOG\,+\,SVM & dlib 68-pt, midpoint & Blink, sequential scan & Movement prohibited & CPM/WPM only \\
Abhaya \emph{et al.}~\cite{ref03} & dlib & dlib 68-pt landmarks & CNN, 9 directions (T9) & No & 15.3\,\% error \\
Rafid \emph{et al.}~\cite{ref04} & MediaPipe & \textbf{MediaPipe iris} & $2\times4$ grid, blink & No & 14.53 CPM \\
Varshini \emph{et al.}~\cite{ref05} & dlib & dlib 68-pt, EAR & Blink, group scan & No & 92.5\,\%, 17 WPM \\
Virtual mouse~\cite{ref06} & MediaPipe & \textbf{MediaPipe Face Mesh iris} & Blink duration & Cited, not implemented & --- \\
Kathpal \emph{et al.}~\cite{ref07} & OpenCV/dlib & Horizontal/vertical gaze ratio & Blink & No & 91\,\% / 85\,\% (glasses) \\
Saraswati \emph{et al.}~\cite{ref08} & Haar cascade & Threshold + integral projection & Pointer, blink & No & 40--50\,cm only \\
Cabatuan \emph{et al.}~\cite{ref09} & --- & Sobel + thinning + Circular Hough & On-screen keyboard & No & FPGA $\equiv$ MATLAB \\
Zhou \emph{et al.}~\cite{ref09b} & --- & SMI ETG 2w (commercial IR) & Adaptive dwell, 600\,ms & Noted, unsolved & --- \\
\midrule
\textbf{This work} & MediaPipe & \textbf{MediaPipe iris, canthus-normalised} & \textbf{Dwell pointing} & \textbf{Implicit (canthus frame)} & $\mathbf{4.37^\circ}$ \textbf{continuous} \\
\bottomrule
\end{tabular}
\end{table*}
% ---------------------------------------------------------------------------


\section{The Vertical-Asymmetry Observation}

The observation that motivates our method is available in the surveyed
literature but is not commented upon there. Waideman and Aquino
Junior~\cite{ref01} report directional accuracy split by axis: $93\,\%$ for
left/right and $76\,\%$ for up/down. The gap is large and the paper does not
explain it.

Consider how the vertical channel is conventionally formed. Let $\mathbf{p}$ be
the iris centre and let $t$, $b$ be the upper and lower eyelid landmarks. The
standard vertical feature is
\begin{equation}
v = \frac{p_y - t_y}{b_y - t_y + \epsilon},
\label{eq:eyelid}
\end{equation}
that is, the iris height expressed as a fraction of the \emph{eyelid aperture}.
The denominator of~\eqref{eq:eyelid} is not a fixed anatomical quantity. It
contracts with every squint, every partial blink, and with ambient brightness.
Any such change alters $v$ without the eye having moved. The horizontal channel,
normalised by the distance between the eye corners, suffers no equivalent
contamination.

An eyelid-normalised vertical channel should therefore be measurably worse than
its horizontal counterpart --- which is what~\cite{ref01} reports. Our method
follows directly: normalise against a structure that does not move.


\section{Front End: Estimating the Coordinate}
\label{sec:frontend}

\subsection{Landmark acquisition}

We use MediaPipe Face Landmarker~\cite{ref10}, which returns 478 landmarks including a
five-point iris ring per eye (indices 468--472 and 473--477) and a
$4\times4$ facial transformation matrix. Detection runs on the GPU where
available and falls back to CPU/WASM otherwise. The medial and lateral canthi
are landmarks 133/33 (left) and 362/263 (right).

\subsection{Canthus-relative normalisation}

Let $\mathbf{i}$, $\mathbf{m}$ and $\mathbf{l}$ denote the iris centre and the
medial and lateral canthi of one eye. We define the eye axis and its
perpendicular,
\begin{equation}
\mathbf{u} = \frac{\mathbf{m} - \mathbf{l}}{\lVert \mathbf{m} - \mathbf{l}\rVert},
\qquad
\mathbf{u}^{\perp} = \begin{bmatrix} -u_y \\ u_x \end{bmatrix},
\end{equation}
and take the iris offset from the canthus midpoint,
$\mathbf{d} = \mathbf{i} - \tfrac{1}{2}(\mathbf{m} + \mathbf{l})$. The two
features per eye are the components of that offset in the eye's own frame,
scaled by the inter-canthal distance $w = \lVert \mathbf{m} - \mathbf{l}\rVert$:
\begin{equation}
c_{\parallel} = \frac{\mathbf{d}\cdot\mathbf{u}}{w},
\qquad
c_{\perp} = \frac{\mathbf{d}\cdot\mathbf{u}^{\perp}}{w}.
\label{eq:canthus}
\end{equation}

Two properties follow. The scale $w$ is a Euclidean distance rather than a
coordinate difference, so it does not collapse as the head turns. And because
$\mathbf{u}$ rotates with the head, \eqref{eq:canthus} is invariant to head roll
\emph{by construction} --- no roll angle is estimated, and none is needed.
Fig.~\ref{fig:geometry} contrasts the two normalisations.

\subsection{Feature vector and calibration}

The full vector concatenates the four conventional ratios, head translation and
scale proxies, a bias term, and the four canthus components
of~\eqref{eq:canthus}, giving $d = 12$. Screen coordinates are obtained by ridge
regression,
\begin{equation}
\mathbf{W} = (\mathbf{F}^{\!\top}\mathbf{F} + \lambda \mathbf{I})^{-1}
             \mathbf{F}^{\!\top}\mathbf{T},
\end{equation}
fitted to a short per-user calibration of nine on-screen targets anchored to key
centres, five fixations each. Samples are weighted by a geometric frame-quality
score; ordinary refinement clicks gathered during use are admitted at a tenth of
the weight of a calibration click, since they are not verified ground truth.

% ---------------------------------------------------------------------------
\begin{figure}[!t]
\centering
\figholder{5.4cm}{%
  \textbf{Fig.~2 placeholder --- normalisation geometry}\\[4pt]
  Left: eyelid-relative. The denominator is the aperture $b_y-t_y$,\\
  which contracts on squint and partial blink.\\[6pt]
  Right: canthus-relative. The offset $\mathbf{d}$ is resolved on\\
  $(\mathbf{u},\mathbf{u}^{\perp})$ and scaled by the inter-canthal distance $w$,\\
  both of which are rigid.\\[8pt]
  \emph{Replace with} \texttt{\textbackslash includegraphics[width=\textbackslash linewidth]\{fig\_geometry\}}}
\caption{Eyelid-relative against canthus-relative normalisation of the iris
offset. Only the latter is invariant to eyelid aperture and to head roll.}
\label{fig:geometry}
\end{figure}
% ---------------------------------------------------------------------------


\subsection{Components integrated around the estimator}
\label{sec:integrated}

The coordinate a regressor emits is not yet a usable pointer, and several
components between the two are load-bearing. We list them because they are the
difference between the numbers in Section~\ref{sec:results} and a system that is
unpleasant to use, and because two of them were added in response to measured
defects rather than to intuition.

\textbf{Temporal filtering.} Predictions pass through a One Euro
filter~\cite{ref11}, whose cutoff rises with speed so that a fixation is smoothed
heavily while a saccade is not lagged. Within the dwell window the committed
point is the fixation centroid rather than the instantaneous estimate.

\textbf{Dwell hysteresis.} A cursor resting on a key boundary would otherwise
oscillate between neighbours and commit neither. The dwell target is retained
until the estimate leaves the key by a fixed margin for several consecutive
frames, which converts boundary jitter into a stable selection.

\textbf{Coverage-aware calibration eviction.} The calibration buffer is bounded.
Evicting oldest-first appears harmless and is not: it silently destroys spatial
coverage, and we measured a buffer that had collapsed from thirteen occupied
grid cells to two, with error rising from $4.93^\circ$ to $8.25^\circ$. Eviction
is therefore driven by grid occupancy, the original calibration anchors are
evicted last, and unverified refinement samples enter at a tenth of the weight
of a calibration click.

\textbf{Geometric quality gating.} Samples are weighted by a frame-quality score
combining face scale, centring and landmark stability, and an adaptive eye
aspect ratio baseline suppresses frames during blinks. The baseline is per-user
because absolute EAR varies more between people than a blink varies within one.

\textbf{User-centred pose convention.} The facial transformation matrix is
right-handed in the camera's frame, which inverts the sense of yaw and pitch
relative to how a user describes their own movement. Uncorrected, a
head-movement recording is labelled backwards. Correcting the signs does not
improve accuracy --- a sign flip is absorbed by the regression weights --- but it
is a prerequisite for any posture-stratified analysis being meaningful, and
Table~\ref{tab:budget} depends on it.

\textbf{Reliability gating.} A movement classifier refuses to commit from
windows that look like saccades or blinks, and a Bayesian network converts its
state estimate into $P(\text{key correct})$, which tempers the fused posterior so
that a marginal frame falls below the commit threshold and asks instead. A
watchdog disables the gate after consecutive refusals, since a gate whose input
drifts out of distribution can otherwise refuse every keystroke.

\textbf{Capture controls.} Exposure and contrast are exposed in the interface,
because the illumination range over which landmark detection is stable is
narrower than the range a webcam will silently auto-expose across.

\section{A Quantum Conditional Selection Rule}
\label{sec:rule}

\subsection{What a product of marginals discards}

At commit time the deployed system scores each candidate key $k$ by multiplying
the gaze posterior with a language prior,
\begin{equation}
\mathrm{score}(k) \;=\; P(k \mid \text{gaze}) \cdot
  \max\!\big(P(k \mid \text{prefix}),\, \varepsilon\big)^{\beta},
\label{eq:loglinear}
\end{equation}
with $\beta < 1$ damping the prior so that the language model advises rather
than vetoes, and $\varepsilon$ flooring it so that no letter becomes untypable.
$P(k \mid \text{gaze})$ is the exact integral of the gaze Gaussian over the key
rectangle; $P(k \mid \text{prefix})$ comes from a frequency-weighted prefix trie
backed off to a character $n$-gram.

The second factor is a \emph{marginal}. It is the total frequency mass of every
vocabulary word that continues the prefix with $k$, and forming that sum
destroys the distinction between a letter supported by one dominant continuation
and one supported by a diffuse spray that agrees on nothing beyond a single
character. After the prefix \texttt{th}, the letters \texttt{e} and \texttt{r}
may carry comparable marginal mass while meaning very different things.
Equation~\eqref{eq:loglinear} cannot tell them apart, because the information
was gone before the product was taken.

\subsection{The rule}

Let $H$ be the vocabulary completions of the current prefix, each with frequency
weight $w_h$ and each proposing a next letter $k_h$. Candidates occupy a
five-qubit register spanning $a$--$z$ and a word-boundary symbol, prepared in
superposition with amplitudes carrying the gaze likelihood:
\begin{equation}
\lvert \psi \rangle \;=\; \sum_{k} \sqrt{P(k \mid \text{gaze})}\; \lvert k \rangle .
\end{equation}

A language-conditioned diagonal unitary then attaches to each candidate both a
magnitude and a phase. Writing $v_h = w_h / P(k \mid \text{prefix})$ for the
hypothesis weights \emph{within} letter $k$, so that $\sum_h v_h = 1$, the
conditioned amplitude is
\begin{equation}
a(k) \;=\; \sqrt{P(k \mid \text{gaze})}\;
           P(k \mid \text{prefix})^{\beta/2}\;
           \underbrace{\sum_{h:\,k_h = k} v_h\, e^{i\varphi_h}}_{c(k)} ,
\label{eq:rule}
\end{equation}
and the Born rule commits the most probable candidate, $P(k) = \lvert a(k)
\rvert^{2}$. The factor $c(k)$ is the \emph{coherence} of letter $k$: a weighted
mean phasor, with
\begin{equation}
\lvert c(k) \rvert^{2}
  = \sum_{h}\sum_{h'} v_h v_{h'} \cos(\varphi_h - \varphi_{h'}) ,
\label{eq:interference}
\end{equation}
whose off-diagonal terms are the interference. Since $\lvert c(k)\rvert \le 1$
with equality exactly when the supporting hypotheses agree in phase, coherence
acts as a penalty on letters reachable only through mutually incoherent
continuations.

Phases are assigned from each hypothesis's mean per-character log-frequency,
rescaled to $[0,1]$ across $H$ and multiplied by a scale $\sigma_\varphi \pi$, so
that the least surprising completion sits at phase $0$. Length must enter:
phase driven by raw frequency alone is degenerate in precisely the case the
mechanism exists to catch, since a spray of comparably frequent completions
would receive one common phase and interfere perfectly constructively.

\subsection{Containment and expressiveness}

\begin{proposition}[Containment]
\label{prop:contain}
At $\sigma_\varphi = 0$, \eqref{eq:rule} reduces to \eqref{eq:loglinear}.
\end{proposition}

\noindent\emph{Proof.} With all $\varphi_h = 0$ we have $c(k) = \sum_h v_h = 1$
for every $k$, so $\lvert a(k)\rvert^{2} = P(k\mid\text{gaze})\,
P(k\mid\text{prefix})^{\beta}$, which is~\eqref{eq:loglinear} up to the shared
normalisation. $\square$

This is not a formality. A model class that cannot express the incumbent cannot
be fairly said to beat it, and the comparison in Section~\ref{sec:results} is
only meaningful because the baseline lies strictly inside the family being
tested. We assert the identity numerically to $10^{-12}$ over randomised inputs
rather than relying on the derivation.

\begin{proposition}[Expressiveness]
\label{prop:express}
No classical mixture over the same hypotheses reproduces~\eqref{eq:interference}.
\end{proposition}

\noindent\emph{Proof sketch.} A mixture assigns letter $k$ the probability
$\sum_h v_h^{2}$ after normalisation --- the diagonal of~\eqref{eq:interference}
alone. The off-diagonal terms depend on phase differences, and a mixture has no
phase. Since $\big(\sum_h v_h\big)^{2} > \sum_h v_h^{2}$ whenever more than one
hypothesis contributes, the two disagree. $\square$

Concretely, two hypothesis sets with \emph{identical} marginals --- so
indistinguishable to~\eqref{eq:loglinear} --- receive different posteriors
under~\eqref{eq:rule}, and the diffuse set is the one penalised.

\subsection{The consistency oracle and amplitude amplification}

The cluster decoder's correctness guarantee is a predicate: a hypothesis is
admissible only if re-encoding it reproduces exactly the cluster sequence the
user gazed. As an operator this is a projector $O_S$ onto the admissible
subspace, and it is what prevents the language model from ever proposing a word
the gaze did not support.

Given a projector, amplitude amplification~\cite{ref13}, the generalisation of
Grover search~\cite{ref14} to a non-uniform initial distribution, rotates
amplitude into the marked subspace in $O(1/\sqrt{p_S})$ rather than $O(1/p_S)$
steps, reflecting about the prepared state:
\begin{equation}
Q \;=\; -\big(2\lvert\psi\rangle\langle\psi\rvert - I\big)\, O_S ,
\qquad
m \;=\; \left\lfloor \frac{\pi/4}{\arcsin\sqrt{p_S}} \right\rfloor .
\end{equation}
The intent is to let the top posterior clear the commit threshold from fewer
frames of accumulated evidence, shortening dwell. Section~\ref{sec:results}
reports what it does instead.

\subsection{Cost}

Twenty-seven candidates occupy a $32$-amplitude register. Preparation,
conditioning, amplification and measurement are a handful of vector operations
costing microseconds, far inside a $33$\,ms frame budget. We claim no quantum
speedup and simulate no hardware: the claim is about the selection rule's
functional form, not its execution. The distinction matters, because the
amplitude encoding used here is a product of per-feature factors and therefore
classically simulable by construction.


\section{Evaluation Protocol}
\label{sec:protocol}

The protocol is described in full because the survey suggests that evaluation
practice, rather than algorithmic choice, is the largest source of
incomparability in this area.

\subsection{Recording}

Each session presents 21 targets: 13 used for fitting and 8 held out at
positions the fit never sees, placed to span the corner, edge and centre regions
of the display. Each target receives five fixations of 300\,ms of retained
samples, with the first 150\,ms discarded to allow the saccade to settle. The
participant's screen dimensions and viewing distance are recorded so that the
conversion to degrees is measured rather than assumed.

\subsection{Splits}

All splits are \textbf{grouped by target}. The five fixations recorded at one
position are one observation measured five times, not five observations;
splitting by sample leaves near-identical siblings in the training set and
selects a regulariser far too small. Unless stated otherwise, results are
leave-one-target-out over all 21 positions.

\subsection{Metrics}

We report three units. \emph{Pixels} are what the dwell hit-test sees.
\emph{Degrees of visual angle} make the work comparable with the gaze
literature. \emph{Fractions of a key pitch} answer whether the error matters for
typing at all: a mean error above one key pitch means the estimate lands on the
wrong key more often than the right one, irrespective of how good the percentage
looks. Throughput figures quoted for comparison follow the conventions
of~\cite{ref12}.

\subsection{Data quality}

Sessions are admitted only if they cover the full target grid and lose less than
$20\,\%$ of frames to tracking dropout. One recorded session lost $67\,\%$ of its
frames to poor illumination and is excluded from all aggregates; we report its
existence rather than silently dropping it.


\subsection{Selection-rule protocol}
\label{sec:ruleproto}

The selection rule is evaluated by end-to-end typing simulation over the
$16{,}839$-word vocabulary, with the gaze likelihood computed as the exact
Gaussian integral over key rectangles --- the same computation the deployed
keyboard performs --- at controlled noise levels $\sigma \in \{30,40,50,60,75\}$
pixels against a $100 \times 110$ pixel key. Three points of protocol carry the
weight of the comparison.

\textbf{Identical noise.} Both arms are driven from the same seeded generator,
so every arm sees the same draws and the same key posteriors. Any difference is
the selection rule and nothing else.

\textbf{Matched prompt rate.} Accuracy alone cannot settle this. The keyboard
declines to commit below a confidence threshold and asks instead, and a prompt
always yields the correct character, so any rule reaches $100\,\%$ accuracy by
prompting constantly. We therefore sweep the threshold, trace the whole
accuracy/prompt-rate curve, and compare \emph{commit accuracy at matched prompt
rate}. This also removes the threshold as a confound, since neither arm keeps a
threshold tuned in its own favour.

\textbf{Held-out phase scale.} $\sigma_\varphi$ is tuned on a $60$-sentence
training split and reported on a disjoint $60$-sentence held-out split. A gain
that does not survive the split is not a gain.

For the dwell measurement, each frame contributes an independent noisy sample
whose evidence is discounted to a quarter of one independent observation.
Consecutive gaze frames are strongly correlated --- the eye is still and the
landmark noise is not white --- and treating them as independent lets any arm
commit within one or two frames, leaving nothing to measure. The discount is
applied identically to both arms.

\section{Results}
\label{sec:results}

\subsection{Feature ablation}

Table~\ref{tab:ablation} gives the principal result. Canthus normalisation
reduces pooled error from $5.97^\circ$ to $4.37^\circ$, a $26.5\,\%$ reduction,
and improves every session individually ($+16.5\,\%$, $+48.0\,\%$, $+11.7\,\%$);
Fig.~\ref{fig:ablation} shows the per-session breakdown.

Two rows isolate the cause. Adding an \emph{iris-ring radius} scale cue --- a
metric distance estimate, since human iris diameter is nearly constant across
adults --- yields $1.5\,\%$, within noise. The gain is therefore attributable
specifically to the canthi and not to having more features. Adding explicit
\emph{head-pose} features on top of the canthus frame yields $-1.0\,\%$: the
canthus frame is already head-attached, so an explicit pose stage re-derives
information the representation has already removed. Head pose alone, without
canthus normalisation, is worth $7.1\,\%$, which is the size of the effect the
canthus frame subsumes.

% ---------------------------------------------------------------------------
\begin{table}[!t]
\renewcommand{\arraystretch}{1.3}
\caption{Feature Ablation. Leave-One-Target-Out Over All 21 Positions, Three
Recorded Sessions. Lower is Better.}
\label{tab:ablation}
\centering
\footnotesize
\begin{tabular}{@{}lccccc@{}}
\toprule
\textbf{Feature set} & \textbf{Dim} & \textbf{S1} & \textbf{S2} &
\textbf{S3}$^{\dagger}$ & \textbf{Pooled} \\
 & & (px) & (px) & (px) & (px / deg) \\
\midrule
Deployed (eyelid-relative)      & 8  & 196.9 & 281.0 & 277.0 & 251.7 / 5.97 \\
\quad + iris-ring radius        & 11 & 194.6 & 274.1 & 275.0 & 247.9 / 5.88 \\
\quad + head pose only          & 14 & 192.4 & 230.8 & 278.3 & 233.8 / 5.55 \\
\textbf{\quad + canthus}        & \textbf{12} & \textbf{164.5} & \textbf{146.2} & \textbf{244.6} & \textbf{185.1 / 4.37} \\
\quad + canthus + radius        & 15 & 164.5 & 146.0 & 242.7 & 184.4 / 4.35 \\
\quad + canthus + pose          & 21 & 159.7 & 148.7 & 250.3 & 186.2 / 4.40 \\
\quad + all of the above        & 24 & 158.2 & 147.5 & 249.4 & 185.0 / 4.37 \\
\bottomrule
\multicolumn{6}{@{}l@{}}{\footnotesize $^{\dagger}$S3 was recorded with deliberate head movement.}
\end{tabular}
\end{table}
% ---------------------------------------------------------------------------

% ---------------------------------------------------------------------------
\begin{figure}[!t]
\centering
\figholder{5.8cm}{%
  \textbf{Fig.~3 placeholder --- per-session ablation}\\[4pt]
  Grouped bars: three sessions $\times$ feature sets,\\
  error in degrees, lower is better.\\[6pt]
  The point the figure must make: canthus wins\\
  \emph{every} session, while iris radius and head pose\\
  do not separate from the baseline.\\[8pt]
  \emph{Replace with} \texttt{\textbackslash includegraphics[width=\textbackslash linewidth]\{fig\_ablation\}}}
\caption{Held-out error by feature set and session. Canthus normalisation is the
only addition that improves all three sessions.}
\label{fig:ablation}
\end{figure}
% ---------------------------------------------------------------------------

\subsection{An error budget that inverts with head pose}

We decompose the held-out error into a \emph{within-fixation} component --- the
scatter of predictions about their own mean during a single fixation, which
averages away with temporal filtering --- and a \emph{between-fixation}
component, the offset of that mean from the true target, which does not. The
deployed system suppresses the former with the speed-adaptive low-pass filter
and dwell-window fixation centroid of Section~\ref{sec:integrated}.

Table~\ref{tab:budget} shows that the balance reverses with head roll. With a
near-still head, noise is $5.3\,\%$ of the budget and the systematic offset
dominates; beyond $10^\circ$ of roll, noise is $57.9\,\%$ and within-fixation
scatter has grown ninefold, from $0.77^\circ$ to $6.85^\circ$.
Fig.~\ref{fig:budget} plots the crossing.

This has a practical consequence and an interpretive one. Practically, it means
that image-quality and temporal-filtering work --- which attacks only the noise
term --- has a ceiling of about $13\,\%$ for a seated, still user and becomes the
dominant lever once the user moves. Interpretively, it suggests that
disagreements in the literature about whether webcam gaze is limited by sensor
noise or by calibration may be disagreements about posture rather than about
systems.

We further isolated where the added noise originates. Per-frame feature scatter
grows $5.5\times$ ($0.116 \rightarrow 0.638$ of the calibration spread) while the
regressor's amplification of it, measured as the leverage
$\mathbf{f}^{\!\top}\mathbf{A}^{-1}\mathbf{f}$, stays flat ($0.02 \rightarrow
0.06$, i.e. $1.01\times \rightarrow 1.03\times$). The instability is therefore in
the landmark detector, not in the fit: the model is not extrapolating, it is
being fed noisier inputs.

% ---------------------------------------------------------------------------
\begin{table}[!t]
\renewcommand{\arraystretch}{1.3}
\caption{Error Budget by Head Roll. The Character of the Error Inverts:
Systematic Offset Dominates a Still Head, Per-Frame Noise Dominates a Tilted
One.}
\label{tab:budget}
\centering
\footnotesize
\begin{tabular}{@{}lcccc@{}}
\toprule
$|\text{roll}|$ & \textbf{Total} & \textbf{Within-fix.} &
\textbf{Between-fix.} & \textbf{Noise share} \\
\midrule
$0$--$2^\circ$    & $3.31^\circ$ & $0.77^\circ$ & $3.26^\circ$ & $5.3\,\%$ \\
$2$--$5^\circ$    & $4.44^\circ$ & $2.08^\circ$ & $3.38^\circ$ & $27.3\,\%$ \\
$\geq 10^\circ$   & $9.20^\circ$ & $\mathbf{6.85^\circ}$ & $5.84^\circ$ & $\mathbf{57.9\,\%}$ \\
\bottomrule
\end{tabular}
\end{table}
% ---------------------------------------------------------------------------

% ---------------------------------------------------------------------------
\begin{figure}[!t]
\centering
\figholder{5.6cm}{%
  \textbf{Fig.~4 placeholder --- error against head roll}\\[4pt]
  Stacked area or twin-line plot: total error, within-\\
  fixation and between-fixation components against\\
  measured $|\text{roll}|$, with the noise share on a\\
  secondary axis crossing $50\,\%$ near $8$--$10^\circ$.\\[8pt]
  \emph{Replace with} \texttt{\textbackslash includegraphics[width=\textbackslash linewidth]\{fig\_budget\}}}
\caption{Decomposition of held-out error against head roll. The crossing point,
where per-frame noise overtakes systematic offset, falls near $8$--$10^\circ$.}
\label{fig:budget}
\end{figure}
% ---------------------------------------------------------------------------

\subsection{Roll invariance}

Rotating a synthetic face through $0$--$40^\circ$ and recomputing the features,
the canthus components of~\eqref{eq:canthus} drift by $\sim$$10^{-15}$, i.e. at
floating-point precision, while an image-axis formulation of the same offset
drifts by $0.069$ --- a substantial fraction of the feature's own range. The
representation remains sensitive to genuine iris movement, so the invariance is
not achieved by the feature going dead.

On recorded data the practical effect is conditional. Stratified by measured
roll, the eye-frame projection is worth $+11.3\,\%$ beyond $14^\circ$ and
$-4.3\,\%$ below $5^\circ$; pooled over all postures it is $+4.5\,\%$. We
therefore present it as a robustness property rather than an accuracy gain, and
note that head tilt remains expensive under either formulation --- error roughly
triples at $16^\circ$ regardless.


\subsection{The selection rule at matched prompt rate}
\label{sec:ruleresult}

Table~\ref{tab:selection} gives the held-out comparison, under the protocol of
Section~\ref{sec:ruleproto}. Tuned on the training
split, the phase scale that maximised commit accuracy was
$\sigma_\varphi = 0.40$, worth $+0.03$ percentage points there. On held-out
sentences that advantage does not appear: across twenty $(\sigma, \text{prompt
rate})$ cells the median change is $-0.03$ percentage points, and a paired
Wilcoxon test over sentences finds no significant difference at any of the five
noise levels --- zero wins and zero losses. Fig.~\ref{fig:selection} plots the
two frontiers, which lie on top of one another.

The nulls are not an artefact of a rule that fails to engage. The two posteriors
do differ, just negligibly and without direction: the mean Kullback--Leibler
divergence between them is $4.2\times10^{-4}$ nats per keystroke. Nor is the
outcome an artefact of a scale chosen too small to matter. Scales above $0.6$
produce substantial interference and \emph{lose}, falling from $97.0\,\%$ to
$95.3\,\%$ on the training split as $\sigma_\varphi$ rises from $0$ to $1.0$.
There is no setting at which the phase channel helps.

\subsection{Amplitude amplification and dwell}

Table~\ref{tab:dwell} reports frames to commit. Two findings, in order of
interest.

First, applied after the language-conditioned unitary, amplification is a
\emph{no-op by construction}. The conditioned state already places $94$--$98\,\%$
of its amplitude on consistent letters, so the optimal round count
$\lfloor(\pi/4)/\arcsin\sqrt{p_S}\rfloor$ evaluates to zero and the operator is
never applied. The quadratic speedup is real, but it addresses a bottleneck this
system does not have --- the admissible subspace is not rare.

Second, forced to act where there \emph{is} scattered amplitude to rotate ---
on the raw gaze register, before conditioning, where the marked mass on a
scattered fixation runs $0.31$--$0.78$ --- amplification makes things worse.
Median dwell is unchanged at $2.0$ frames, mean dwell rises from $2.63$ to
$2.80$, and commit accuracy falls from $98.4\,\%$ to $96.4\,\%$. A paired test
over the $249$ keystrokes both arms committed gives $p = 0.48$. We therefore do
not support the claim that amplitude amplification shortens dwell time in this
system; on our measurements it lengthens it slightly while costing accuracy.

\subsection{Why interference fails: a sign, not a magnitude}
\label{sec:why}

A null result is more useful with a mechanism attached, and this one has a sharp
one. For interference to help, the coherence penalty must fall more heavily on
wrong letters than on right ones. Table~\ref{tab:coherence} measures exactly
that separation, and it is negative at every phase scale and grows steadily more
negative as the scale rises, from $-0.004$ at $\sigma_\varphi = 0.2$ to
$-0.108$ at $\sigma_\varphi = 2.0$.

The cause is structural. The correct letter following an English prefix is
typically the one with the \emph{richest} vocabulary support: many genuine
continuations, therefore many hypotheses, therefore the widest spread of phases
and the lowest coherence. A wrong letter is typically backed by one obscure
completion and sits at coherence $1$. Any phase computed from hypothesis
likelihood thus penalises the letter with the strongest linguistic evidence.

We checked two alternative explanations and rejected both. The effect is not a
shortage of hypotheses: raising the budget from $4$ to $4096$ completions moves
the mean coherence only from $0.982$ to $0.988$, whereas equal-weight
hypotheses with phases spread over $[0,\pi]$ would sit between $0.33$ and
$0.64$. Nor is it Zipfian weight concentration:
flattening the within-letter weights to uniform, holding phases fixed, changes
the mean coherence from $0.9883$ to $0.9879$. The limitation is the phase
\emph{source}. A language model offers no feature orthogonal to the likelihood
that already determines the magnitude, so the phase channel carries no
independent information for interference to exploit.

We state the scope of this negative result carefully. It concerns
likelihood-derived phase on a letter-prediction task, not the selection rule as
a formal object, and not tasks where an orthogonal phase source exists.
Propositions~\ref{prop:contain} and~\ref{prop:express} are unaffected.

% ---------------------------------------------------------------------------
\begin{table}[!t]
\renewcommand{\arraystretch}{1.3}
\caption{Selection Rule Against Deployed Log-Linear Fusion, Held-Out, at Matched
Prompt Rate. Each Cell Gives the Baseline Commit Accuracy and the Rule's Change.
$p$ is a Paired Wilcoxon Over Sentences at the Deployed Threshold.}
\label{tab:selection}
\centering
\footnotesize
\begin{tabular}{@{}lcccc c@{}}
\toprule
$\sigma$ (px) & \textbf{2\,\% prompts} & \textbf{5\,\%} & \textbf{10\,\%} &
\textbf{20\,\%} & $p$ \\
\midrule
30 & 0.971 \; $+$0.003 & 0.977 \; $+$0.002 & 0.986 \; $-$0.001 & 0.988 \; $-$0.000 & 1.00 \\
40 & 0.921 \; $-$0.002 & 0.932 \; $+$0.001 & 0.958 \; $-$0.000 & 0.973 \; $+$0.003 & 0.80 \\
50 & 0.888 \; $-$0.002 & 0.897 \; $-$0.001 & 0.914 \; $-$0.004 & 0.945 \; $-$0.001 & 0.34 \\
60 & 0.841 \; $-$0.000 & 0.853 \; $-$0.000 & 0.862 \; $+$0.001 & 0.891 \; $+$0.008 & 0.80 \\
75 & 0.778 \; $-$0.000 & 0.795 \; $-$0.003 & 0.819 \; $+$0.001 & 0.851 \; $-$0.004 & 0.80 \\
\midrule
\multicolumn{6}{@{}l@{}}{Median change over all 20 cells: $\mathbf{-0.0003}$.
Significant wins: \textbf{0}. Losses: \textbf{0}.} \\
\bottomrule
\end{tabular}
\end{table}
% ---------------------------------------------------------------------------

% ---------------------------------------------------------------------------
\begin{table}[!t]
\renewcommand{\arraystretch}{1.3}
\caption{Frames to Commit, Held-Out, $\sigma = 40$\,px. Amplification Does Not
Shorten Dwell; Forced to Act, It Lengthens the Mean and Costs Accuracy.}
\label{tab:dwell}
\centering
\footnotesize
\begin{tabular}{@{}lccc@{}}
\toprule
\textbf{Rule} & \textbf{Median} & \textbf{Mean} & \textbf{Accuracy} \\
\midrule
Log-linear (deployed)            & 2.0 & \textbf{2.63} & \textbf{0.984} \\
Interference only                & 2.0 & 2.66 & 0.988 \\
Amplification, pre-conditioning  & 2.0 & 2.80 & 0.964 \\
Amplification, post-conditioning$^{\ddagger}$ & 2.0 & 2.80 & 0.968 \\
Interference $+$ amplification   & 2.0 & 2.75 & 0.964 \\
\bottomrule
\multicolumn{4}{@{}p{7.6cm}@{}}{\footnotesize $^{\ddagger}$Zero rounds fire:
after conditioning the admissible subspace already holds $94$--$98\,\%$ of the
amplitude. Paired over the 249 keystrokes committed by both arms, $p = 0.48$.}
\end{tabular}
\end{table}
% ---------------------------------------------------------------------------

% ---------------------------------------------------------------------------
\begin{table}[!t]
\renewcommand{\arraystretch}{1.3}
\caption{Why Interference Fails. Coherence on the Correct Letter Against the
Wrong Ones. A \emph{Positive} Separation Would Mean Interference Protects the
Correct Letter; It Is Negative Everywhere and Worsens With Scale.}
\label{tab:coherence}
\centering
\footnotesize
\begin{tabular}{@{}cccc@{}}
\toprule
$\sigma_\varphi$ & $\lvert c \rvert$ \textbf{correct} &
$\lvert c \rvert$ \textbf{wrong} & \textbf{Separation} \\
\midrule
0.2 & 0.9934 & 0.9971 & $-0.0038$ \\
0.4$^{*}$ & 0.9740 & 0.9887 & $-0.0147$ \\
0.8 & 0.9046 & 0.9572 & $-0.0526$ \\
1.0 & 0.8607 & 0.9357 & $-0.0750$ \\
2.0 & 0.7691 & 0.8768 & $\mathbf{-0.1076}$ \\
\bottomrule
\multicolumn{4}{@{}l@{}}{\footnotesize $^{*}$Scale selected on the training split.}
\end{tabular}
\end{table}
% ---------------------------------------------------------------------------

% ---------------------------------------------------------------------------
\begin{figure}[!t]
\centering
\figholder{5.6cm}{%
  \textbf{Fig.~5 placeholder --- the accuracy/prompt-rate frontier}\\[4pt]
  Commit accuracy against prompt rate, one curve per arm,\\
  at $\sigma = 40$ and $\sigma = 60$\,px. The two curves should\\
  overlie each other almost exactly --- that \emph{is} the result.\\[6pt]
  Inset: coherence separation against phase scale, crossing\\
  zero nowhere and falling monotonically.\\[8pt]
  \emph{Replace with} \texttt{\textbackslash includegraphics[width=\textbackslash linewidth]\{fig\_selection\}}}
\caption{The selection rule traces the same accuracy/prompt-rate frontier as the
fusion it generalises. The inset gives the reason: the coherence penalty is
anti-correlated with correctness at every phase scale.}
\label{fig:selection}
\end{figure}
% ---------------------------------------------------------------------------


\section{Negative Results}

We record the following because each was implemented and measured, and because
the cost of re-attempting them is otherwise borne by the next group. The first
two are the subject of Section~\ref{sec:ruleresult} and are restated here only
so the list is complete.

\textbf{Phase interference between language hypotheses.} Median $-0.03$
percentage points at matched prompt rate; no significant win or loss at any of
five noise levels. Retained as a research module, not deployed. The mechanism of
the failure is a sign rather than a magnitude (Section~\ref{sec:why}).

\textbf{Amplitude amplification on the consistency oracle.} Zero rounds fire
after conditioning, because the admissible subspace is not rare. Forced to act,
mean dwell rises from $2.63$ to $2.80$ frames and commit accuracy falls from
$98.4\,\%$ to $96.4\,\%$.

\textbf{Quantum-kernel regression.} A fidelity kernel
$K(x,z)=|\langle\phi(x)|\phi(z)\rangle|^2$ was evaluated in place of the linear
model, using both angle and ZZ~\cite{ref15} feature maps at 12 qubits. The angle map
scored $-104\,\%$ against the linear baseline and the ZZ map $-989\,\%$. The
failure of the latter is the expected exponential
concentration~\cite{ref16}: at 12 qubits the fidelity
between deliberately similar samples falls to $0.005$, the Gram matrix approaches
the identity, and the model can only memorise. A classical Gaussian kernel on the
same features scored $+9.7\,\%$. We note additionally that the angle-map kernel
is a product of per-feature cosines and therefore classically simulable; no
quantum advantage is claimed or observed.

\textbf{Nonlinear regression generally.} A tuned RBF kernel ridge regressor beat
the linear model by $20.2\,\%$ on the eyelid-relative features but by only
$0.5\,\%$ once canthus normalisation was present. The curvature the kernel was
modelling was an artefact of the contaminated feature, not a property of gaze
geometry.

\textbf{Discrete hierarchical selection.} Two-level region selection is
$\sim$20 percentage points more reliable per choice than pointing directly at a
key ($55.0$/$68.8\,\%$ against $35.2$/$46.5\,\%$), but doubling the dwells costs
more than the reliability returns: throughput falls from $18.5$ to $14.4$ and
from $24.4$ to $18.1$ correct letters per minute.

\textbf{Eye-region de-rotation.} Given that landmark jitter grows $5.5\times$
under roll, de-rotating the frame before detection --- presenting the detector
with an upright face --- is an obvious remedy. Implemented with an exact inverse
mapping, it produced no perceptible gain in live use. We suspect the detector's
internal face alignment already performs this normalisation.


\section{Limitations}
\label{sec:limits}

The results rest on \textbf{three recorded sessions of a single participant}.
This is sufficient to establish a direction --- the canthus effect is large and
consistent in sign across all three --- but not to establish an effect size that
will generalise, and we do not report significance testing. Inter-personal
variation in eye geometry is precisely what a canthus-relative frame is supposed
to be robust to, which makes a multi-participant study the natural next step and
the most informative one.

The mechanism argument, that eyelid aperture contaminates the vertical channel,
is supported by the ablation and by the asymmetry reported in~\cite{ref01}, but
is not yet demonstrated causally. The decisive experiment is a session in which
eyelid aperture is deliberately perturbed --- prompted squinting --- while gaze
target is held fixed. The prediction is specific: eyelid-relative features should
degrade sharply and canthus-relative features should not.

The selection-rule results carry a different set of caveats. They come from
simulation rather than from live typing: the gaze likelihood is the deployed
key-rectangle integral driven by Gaussian noise at controlled $\sigma$, which
models isotropic error and not the anisotropic, temporally correlated error a
real tracker produces. The dwell measurement discounts each frame to a quarter
of an independent observation to stand in for that correlation, but the discount
is a stipulation rather than a measurement. The negative result is also specific
to one vocabulary and one language: our claim is that likelihood-derived phase
is mis-signed for English letter prediction, not that interference-based fusion
is useless in general. A task supplying a phase source genuinely orthogonal to
the likelihood --- multimodal evidence, or a syntactic signal independent of
frequency --- would be a fair test of the mechanism and we have not run one.

Finally, all sessions were recorded at one display and one nominal viewing
distance. The degrees conversion is measured rather than assumed, but the
generalisation across geometries is untested.


\section{Conclusion}

We set out to improve how a gaze keyboard decides, and improved how it estimates
instead.

The quantum conditional selection rule is a strict generalisation of the
log-linear fusion that gaze keyboards conventionally use. It places candidates
in superposition with gaze-derived amplitudes, conditions them on language
through both magnitude and phase, constrains them with a consistency-preserving
oracle, and commits by the Born rule. It recovers the incumbent exactly at zero
phase, and its interference cross-terms are provably beyond any classical
mixture of the same hypotheses. Measured against that incumbent on identical
noise at matched prompt rate, it is worth a median $-0.03$ percentage points,
and amplitude amplification lengthens dwell rather than shortening it.

We regard the explanation as the more durable contribution. Interference fails
here for a reason that no retuning addresses: the correct letter after an
English prefix is the one with the richest vocabulary support, hence the most
hypotheses, the widest phase spread and the lowest coherence. The measured
separation between the coherence of correct and incorrect letters is negative at
every phase scale and worsens monotonically with it. Phase derived from
hypothesis likelihood penalises exactly the evidence it should reward. For
interference to be useful, phase must come from a source orthogonal to the
magnitude --- and a language model does not provide one.

The accuracy came from the front end. Measuring the continuous coordinate in
degrees, rather than as region classification, surfaced a defect in a step the
literature treats as settled: normalising the iris offset by the eyelid aperture
couples the vertical gaze channel to a structure that moves independently of
gaze. Normalising against the canthi instead reduces held-out error by
$26.5\,\%$, wins on every recorded session, explains an axis asymmetry reported
but undiagnosed in prior work, and makes an explicit head-pose stage
unnecessary. The change is four features and no additional runtime cost. We also
find that the error budget inverts with posture --- from $5.3\,\%$ per-frame
noise with a still head to $57.9\,\%$ beyond $10^\circ$ of roll --- and that the
added instability originates in the landmark detector rather than the regressor.

Immediate future work is the eyelid-perturbation experiment described in
Section~\ref{sec:limits}, which would convert the canthus mechanism from a
supported inference into a demonstrated one, and a multi-participant study to
establish whether the effect size holds across eye geometries. For the selection
rule, the open question is whether a phase source orthogonal to likelihood
exists for this task; until one does, we report the rule as a formal
generalisation whose empirical value is nil, and we would rather say so than
leave the next group to discover it.


\section*{Acknowledgment}

The authors thank \dots

\vspace{4pt}
\noindent\textbf{Reproducibility.} The evaluation harness, the recording
protocol, and the per-session results are available with the implementation.


\begin{thebibliography}{10}
\footnotesize

\bibitem{ref01}
B.~Waideman and P.~T. Aquino Junior, ``Augmentative and alternative
communication using eye tracking and word recommendation using language
models,'' \emph{IEEE Latin America Transactions}, vol.~23, no.~8, pp.~637--646,
Aug. 2025, doi: 10.1109/TLA.2025.11072932.

\bibitem{ref02}
M.~R. Islam, M.~S. Rahman \emph{et al.}, ``Computer vision based eye gaze
controlled virtual keyboard for people with quadriplegia,'' in \emph{Proc. 2021
Int. Conf. Automation, Control and Mechatronics for Industry 4.0 (ACMI)},
Rajshahi, Bangladesh, Jul. 2021, doi: 10.1109/ACMI53878.2021.9528213.

\bibitem{ref03}
V.~Abhaya, A.~S. Bharadwaj, C.~C. Bagan, K.~Dhanraj, and G.~Shyamala,
``Eye-Move: An eye gaze typing application with OpenCV and Dlib library,'' in
\emph{Proc. 2022 Int. Conf. Automation, Computing and Renewable Systems
(ICACRS)}, 2022, doi: 10.1109/ICACRS55517.2022.10029276.

\bibitem{ref04}
S.~T.~S. Rafid, A.~J. Ajwad, and A.~K.~M.~A.~M. Azad, ``Empowering quadriplegic
patients: A computer vision based approach for keyboard control using eye gaze
and blink,'' in \emph{Proc. 2023 26th Int. Conf. Computer and Information
Technology (ICCIT)}, Cox's Bazar, Bangladesh, Dec. 2023,
doi: 10.1109/ICCIT60459.2023.10441390.

\bibitem{ref05}
B.~G. Varshini, A.~Manimegalai, and S.~D. Abirami, ``Eye-gaze controlled virtual
keyboard using blink detection for accessible hands-free communication,'' in
\emph{Proc. 2025 Int. Conf. Computing and Communication Technologies
(ICCCIT)}, 2025, doi: 10.1109/ICCCIT62592.2025.10927769.

\bibitem{ref06}
``Eye gaze tracking for virtual mouse control using computer vision,'' in
\emph{Proc. 2026 4th Int. Conf. Artificial Intelligence and Machine Learning
Applications (AIMLA)}, 2026, doi: 10.1109/AIMLA67915.2026.11522321.

\bibitem{ref07}
K.~Kathpal, S.~Negi, and S.~Sharma, ``iChat: Interactive eyes for specially
challenged people using OpenCV Python,'' in \emph{Proc. 2021 9th Int. Conf.
Reliability, Infocom Technologies and Optimization (ICRITO)}, Noida, India,
2021, doi: 10.1109/ICRITO51393.2021.9596502.

\bibitem{ref08}
V.~I. Saraswati, R.~Sigit, and T.~Harsono, ``Eye gaze system to operate virtual
keyboard,'' in \emph{Proc. Int. Electronics Symp.}, Electronics Engineering
Polytechnic Institute of Surabaya, Indonesia, ISBN 978-1-5090-1640-2.

\bibitem{ref09}
M.~Cabatuan \emph{et al.}, ``Implementation of eye gaze tracking technique on
FPGA-based on-screen keyboard system using Verilog and MATLAB,'' in \emph{Proc.
2017 IEEE Region 10 Conf. (TENCON)}, Malaysia, Nov. 2017, pp.~2772--2777.

\bibitem{ref09b}
X.~Zhou \emph{et al.}, ``Dwell time adaptive three-way text entry by gaze,'' in
\emph{Proc. 31st Chinese Control and Decision Conf. (CCDC)}, Nanjing, China,
2019, pp.~5297--5301.

\bibitem{ref10}
C.~Lugaresi \emph{et al.}, ``MediaPipe: A framework for building perception
pipelines,'' \emph{arXiv:1906.08172}, 2019.

\bibitem{ref11}
G.~Casiez, N.~Roussel, and D.~Vogel, ``The one Euro filter: A simple speed-based
low-pass filter for noisy input in interactive systems,'' in \emph{Proc. ACM
CHI}, 2012, pp.~2527--2530.

\bibitem{ref12}
T.~Soukoreff and I.~S. MacKenzie, ``Metrics for text entry research,'' in
\emph{Proc. ACM CHI}, 2003, pp.~113--120.

\bibitem{ref13}
G.~Brassard, P.~H{\o}yer, M.~Mosca, and A.~Tapp, ``Quantum amplitude
amplification and estimation,'' in \emph{Quantum Computation and Information},
AMS Contemporary Mathematics, vol.~305, 2002, pp.~53--74.

\bibitem{ref14}
L.~K. Grover, ``A fast quantum mechanical algorithm for database search,'' in
\emph{Proc. 28th Annu. ACM Symp. Theory of Computing (STOC)}, 1996,
pp.~212--219.

\bibitem{ref15}
V.~Havl\'{\i}\v{c}ek, A.~D. C\'orcoles, K.~Temme, A.~W. Harrow, A.~Kandala,
J.~M. Chow, and J.~M. Gambetta, ``Supervised learning with quantum-enhanced
feature spaces,'' \emph{Nature}, vol.~567, pp.~209--212, 2019.

\bibitem{ref16}
S.~Thanasilp, S.~Wang, M.~Cerezo, and Z.~Holmes, ``Exponential concentration in
quantum kernel methods,'' \emph{Nature Communications}, vol.~15, art.~5200,
2024.

\bibitem{ref17}
P.~O. Kristensson and S.~Zhai, ``SHARK$^2$: A large vocabulary shorthand writing
system for pen-based computers,'' in \emph{Proc. ACM UIST}, 2004, pp.~43--52.

\end{thebibliography}

\end{document}
